\scanpage{51}
\mspage{50}
\mspage{73}
\margin{n°382}
s'identifie à $\mathrm{Hom}^L(\bigotimes_{i=1}^{i=n}E_i,E)$, cf. n° 6.7.5. Soit U un ouvert de B, et soient $s_i\in S_{E_i}(U)$, $\varphi\in S_F(U)$ ; alors $\otimes s_i$ est une section $\bigotimes_{i=1}^{i=n}E_i$, et comme $\varphi$ s'identifie à un morphisme de $\bigotimes_{i=1}^{i=n}E_i$ dans E, $\varphi$ transforme $\otimes s_i$ en une section de E ; on la note $\langle\varphi,s_1\otimes\cdots\otimes s_n\rangle$, ou encore $\varphi(s_1,\ldots,s_n)$. C'est une fonction $C_L(U)$-multilinéaire de $\varphi$ et des $s_i$.

Inversement, supposons donné, pour tout ouvert U de B, une \uline{application} $C_L(U)$-\uline{multilinéaire} $\varphi_U$ \uline{de} $S_{E_1}(U)\times\cdots\times S_{E_n}(U)$ \uline{dans} $S_E(U)$, le système $(\varphi_U)$ \uline{commutant aux homomorphismes de restriction} (cf. n° 6.2.3. pour le cas $n=1$). Il existe alors une section $\varphi$ de F sur B et une seule qui correspond aux $\varphi_U$.

Lorsque tous les $E_i$ sont égaux, et que l'on remplace F par $F'=\bigwedge^n E_1^*\otimes E=\mathrm{Hom}^L(\bigwedge^n E_1,E)$ (cf. n° 6.8.7), on a un résultat analogue, les applications $\varphi_U$ étant cette fois des applications multilinéaires \uline{alternées}.

6.8 -- \uline{Espaces tensoriels}, \uline{puissances extérieures}, \uline{puissances symétriques}.

Les espaces fibrés vectoriels considérés dans ce numéro sont relatifs à un corps L et à une base B fixés ; on suppose que L est \uline{commutatif}.

6.8.1. Soit E un espace fibré vectoriel, et soient I et J deux ensembles finis disjoints. Le \uline{fibré tensoriel} $T_J^I(E)$ est défini comme le produit tensoriel $\bigotimes_{\alpha\in I\cup J}E_\alpha$, où $E_\alpha=E$ si $\alpha\in I$, et $E_\alpha=E^*$ si $\alpha\in J$, $E^*$ désignant le dual de E (cf. \uline{Alg.}, chap. III, 3ème édition, \S\ \qquad, n°\qquad). La fibre $T_J^I(E)_b$ de ce fibré en un point b est égale à l'espace tensoriel $T_J^I(E_b)$ défini en \uline{Alg.}, \uline{loc.} \uline{cit.}. Lorsque $I=\{1,\ldots,p\}$ et $J=\{p+1,\ldots,p+q\}$, on écrit $T_q^p(E)$ au lieu de $T_J^I(E)$ ; on a
\[
T_q^p(E)=(\otimes^p E)\otimes(\otimes^q E^*).
\]
La donnée d'un \uline{ordre total} sur I et sur J définit un isomorphisme canonique de $T_J^I(E)$ sur $T_q^p(E)$.
